% Independent preview; original geometry and numeric relationships retained.
% The four local pairs occur in the support, but their joint configuration does not.
\begin{tikzpicture}[
  x=1mm,y=1mm,font=\figurefont\footnotesize,
  state/.style={circle,draw=NNHidden,fill=NNHidden!18,
    minimum size=8mm,inner sep=0pt,line width=1.1pt},
  edge/.style={draw=NNWire,line width=1.2pt},
  pair/.style={font=\figurefont\footnotesize,text=FigureInk,inner sep=1mm},
  index/.style={font=\figurefont\footnotesize,text=FigureMuted},
  note/.style={align=left,text=FigureInk,anchor=west}
]
  \node[state] (x1) at (16,34) {$0$};
  \node[state] (x2) at (46,34) {$0$};
  \node[state] (x3) at (46,4) {$1$};
  \node[state] (x4) at (16,4) {$0$};
  \draw[edge] (x1)--node[pair,above] {$(0,0)$} (x2);
  \draw[edge] (x2)--node[pair,right] {$(0,1)$} (x3);
  \draw[edge] (x4)--node[pair,below] {$(1,0)$} (x3);
  \draw[edge] (x1)--node[pair,left] {$(0,0)$} (x4);
  \node[index,anchor=south east] at (12,38) {$x_1$};
  \node[index,anchor=south west] at (50,38) {$x_2$};
  \node[index,anchor=north west] at (50,0) {$x_3$};
  \node[index,anchor=north east] at (12,0) {$x_4$};
  \node at (31,19) {$\bm y$};
  \node[font=\figurefont\footnotesize] at (86,38)
    {支持集：$U,V$的$2\times2$表};
  \node at (76,31) {$V=0$};
  \node at (100,31) {$V=1$};
  \foreach \y/\label/\left/\right in {
    20/{U=0}/{0000}/{0110},
    6/{U=1}/{1011}/{1101}}{
    \node[anchor=east] at (62,\y) {$\label$};
    \foreach \x/\value in {64/\left,88/\right}{
      \fill[FigureBlueFill] (\x,{\y-6}) rectangle ++(24,12);
      \draw[FigureStroke,line width=.6pt] (\x,{\y-6}) rectangle ++(24,12);
      \node[align=center] at ({\x+12},\y) {$\value$\\概率$1/4$};
    }
  }
  \node[draw=FigureRed,fill=FigureRedFill,rounded corners=.7mm,
    line width=.8pt,align=center,inner sep=2pt] at (86,-9)
    {其余12个配置概率为0\\特别地，$p(0010)=0$};
\end{tikzpicture}
